\ifx\allfiles\undefined
\documentclass[12pt, a4paper, oneside, UTF8]{ctexbook}
\def\path{../config}
\input{../config/_config}
\begin{document}

% \input{../config/cover}
\else
\fi
\chapter{真题与模拟题}

\begin{enumerate}
    \item [(91)]设4阶矩阵$A=\begin{pmatrix}
    5 & 2 & 0 & 0 \\
    2 & 1 & 0 & 0 \\
    0 & 0 & 1 & -2 \\
    0 & 0 & 1 & 1
    \end{pmatrix}$ 则A的逆矩阵$A^{-1}$为

    \item [(91)]设n是曲面$2x^2+3y^2+z^2=6$在点$P(1,1,1)$处的指向外侧的法向量,求函数$u=\frac{\sqrt{6x^2+8y^2}}{z}$在
    点P处的沿方向n的方向导数是?

    \item[(91)] \bl 将函数$f(x)=2+\left|x\right|(-1\leq x\leq 1)$ 展开成以2为周期的傅里叶级数,并求此级数$\sum_{n=1}^{\infty}\frac{1}{n^2}$
    的和. 

    \item[(91)] \bl 在上半平面求一条向上凹的曲线,其上任意一点P(x,y)处的曲率等于此曲线在该点法线段PQ长度的倒数(Q是法线与x轴的交点)
    且曲线在点(1,1)处的切线与x轴平行 

    \item[(91-2)] \bl 求$\int_{3}^{+\infty}\frac{\d x}{(x-1)^4\sqrt{x^2-2x}}$
    \item[(91-2)] \bt 计算$I=\iint_{S}-y\d z\d x + (z+1)\d x\d y$其中S是圆柱面$x^2+y^2=4$被平面$x+z=2$和$z=0$ 
    所截出部分的外侧.

    \item[(92)] 在曲线$x=t, y = -t^2, z = t^3$的所有切线中,与平面$x+2y+z=4$平行的切线有(\qquad)条
    \item[(92)] 计算曲面积分$\iint_{\sum}(x^3+az^2)\d y\d z + (y^3 + ax^2)\d z\d x + (z^3 + ay^2)\d x\d y$, 其中$\sum$为上半球面
    $z=\sqrt{a^2 - x^2 - y^2}$的上侧. 

    \item [(92)] 设$f''(x) < 0, f(0) = 0$ 证明对于任意$x_1 > 0, x_2 > 0$ 有$f(x_1+x_2) < f(x_1) + f(x_2)$ 
    \item [(93)] 由曲线$\begin{cases}
        3x^2 + 2y^2 = 12 \\
        z = 0
    \end{cases}$ 绕$y$轴旋转一周得到的旋转曲面在点($0, \sqrt{3}, \sqrt{2}$)处指向外侧的单位法向量为\_\_\_\_
    \item [(93)]设数量场$u=\ln\sqrt{x^2+y^2+z^2}$则$div(grad\ u) = \_\_\_\_$
    \item [(93)] \bt 求级数$\sum_{n=0}^{\infty}\frac{(-1)^n(n^2-n+1)}{2^n}$的和
    \item [(93)]
    \begin{enumerate}
        \item [(1)] 设在[0, $\infty$]上函数f(x)有连续导数, 且$f'(x)\geq k\geq 0, f(0) < 0$ 证明$f(x)$在($0,\infty$)内
        至少有一个零点. 
        \item [(2)] 设$b>a>e$证明$a^b > b^a$ 
    \end{enumerate}
    \item [(93)] 已知$R^3$的两个基为 
    $$
    \begin{cases}
        \alpha_1 = (1,1,1)^T \\
        \alpha_2 = (1,0,-1)^T \\
        \alpha_3 = (1,0,1)^T
    \end{cases} \qquad
    \begin{cases}
        \beta_1 = (1,2,1)^T \\
        \beta_2 = (2,3,4)^T \\
        \beta_3 = (3,4,3)^T
    \end{cases}
    $$
    求由基$\alpha_i$到基$\beta_i$的过度矩阵P.

    \item [(94)] 曲面$z-e^x+2xy=3$在点($1,2,0$)处的切平面的方程为\_\_\_\_
    \item [(94)] 设$u=e^{-x}\sin\frac{x}{y}$则$\frac{\P^2 u}{\P x\P y}$在点$(2,\frac{1}{\pi})$处的值为 
    \item [(94)] 将函数$f(x)=\frac{1}{4}\ln\frac{1+x}{1-x}+\frac{1}{2}\arctan{x}-x$展开成$x$的幂级数
    \item [(94)] 求$\int \frac{\d x}{\sin {2x} + 2\sin{x}}$
    \item [(94)] 计算曲面积分$\iint_{s}\frac{x\d y\d z + z^2\d x\d y}{x^2+y^2+z^2}$其中$S$是由曲面$x^2+y^2=R^2$及两平面
    $z=R,z=-R(R>0)$锁围成体表面的外侧
    \item [(94)] 设$f(x)$具有二阶连续导数, $f(0)=0, f'(0)=1,$且$[xy(x+y)-f(x)y]\d x + [f'(x)+x^2y]\d y = 0$
    为一阶全微分方程,求$f(x)$及此全微分方程的通解.
    \item [(94)] 已知点A和点B的直角坐标分别为$(1,0,0)$与$(0,1,1)$线段AB绕z轴旋转一周所围成的旋转曲面S,求由S和两平面
    $z=0, z= 1$所围成立体的体积.
    
    \item [(95)] 设$(a\times b)\cdot c = 2$则$[(a+b)\times (b+c)]\cdot(c+a)$=
    \item [(95)] 幂级数$\int_{n=1}^{\infty}\frac{n}{2^n + (-3)^n}x^{2n-1}$的收敛半径是
    \item [(95)] 设三阶方阵$A,B$满足关系是$A^{-1}BA=6A+BA$其中$A=\begin{pmatrix}
        \frac{1}{3} & 0 & 0 \\
        0 & \frac{1}{4} & 0 \\
        0 & 0 & \frac{1}{7}
    \end{pmatrix}$则B = 
    \item [(95)] 计算曲面积分$\iint_{\sum}z\d S$其中$\sum$为锥面$z=\sqrt{x^2+y^2}$在柱体$x^2+y^2\leq 2x$内的部分
    \item [(95)] 将函数$f(x)=x-1(0\leq x\leq 2)$展开成周期为4的余弦函数
    \item [(95)] 设曲线L位于$xoy$平面的第一象限内, L上任意一点M处的切线与y轴总相交, 交点及为A, 已知$\left|MA\right| = \left|OA\right|$
    且L过点($\frac{3}{2},\frac{3}{2}$)求L的方程.
    \item [(95)] 设函数$Q(x,y)$在$xoy$平面上具有一阶连续偏导数,曲线积分$\int_{L}2xy\d x + Q(x,y)\d y$与路径无关,并且对于任意t有
    $$
    \int_{(0,0)}^{(t, 1)}2xy\d x + Q(x,y)\d y = \int_{(0,0)}^{(1,t)}2xy\d x + Q(x,y)\d y
    $$
    求$Q(x,y)$
    \item [(95)] 函数$f(x)$和$g(x)$在$[a,b]$上存在二阶导数,并且$g''(x)\neq 0, f(a) = f(b) = g(a) = g(b) = 0$证明:
    \begin{enumerate}
        \item [(1)] 在开区间(a,b)内$g(x)\neq 0$
        \item [(2)] 在开区间(a,b)内至少存在一点$\xi$, 使得$\frac{f(\xi)}{g(\xi)}=\frac{f''(\xi)}{g''(\xi)}$
    \end{enumerate}
    \item [(95-2)] 计算二重积分$\iint_{D}x^2y\d x\d y,$其中D是由双曲线$x^2-y^2=1$及直线$y=0,y=1$所围成区域

    \item [(96)] 设$f(x)$有二阶连续导数, 且$f'(0)=0,\lim_{x\to 0}\frac{f''(x)}{\left|x\right|} = 1$则
    \begin{choices}[1]
        \task $f(0)$是$f(x)$的极小值
        \task $f(0)$是$f(x)$的极大值
        \task $(0, f(0))$是$f(x)$的拐点 
        \task $f(0)$不是$f(x)$极值,$(0,f(0))$也不是其拐点
    \end{choices}
    \item [(96)] 设$f(x)$有连续导数, $f(0) = 0, f'(0) \neq 0, F(x) = \int_{0}^{x}(x^2-t^2)f(t)\d t$且当$x\to 0, F'(x)$与
    $x^k$时同阶无穷小,则$k=$ 
    \item [(96)] 求心型线$r=a(1+\cos\theta)$的全长,其中$a>0$是常数
    \item [(96)] 计算曲面积分$\iint_S(2x+z)\d y\d z + z\d \d y$其中S为有向曲面$z=x^2+y^2(0\leq z\leq 1)$,其法向量与z轴正向的夹角为锐角.
    \item [(96)] 设交换$\begin{cases}
        u = x - 2y \\
        v = x + ay 
    \end{cases}$ 可以把方程$6\frac{\P^2 z}{\P x^2} + \frac{\P^2 z}{\P x\P y}-\frac{\P^2 z}{\P^2 y}=0$简化为$\frac{\P^2 z}{\P u\P v} = 0$ 
    求常数a.
    \item [(96)] 求级数$\sum_{n=2}^{\infty}\frac{1}{n^2-1}2^n$的和
    \item [(96)] 设$f(x)$在[0,1]上具有二阶导数,且满足条件$\left|f(x)\right|\leq a, \left|f''(x)\right|\leq b$其中a,b均为非负常数,c是
    (0,1)上任意一点.
    \begin{enumerate}
        \item [(1)] 写成$f(x)$在点$x=c$处带拉格朗日型余项的一阶泰勒公式. 
        \item [(2)] 证明$\left|f'(x)\right|\leq 2a+\frac{b}{2}$
    \end{enumerate}
    \item [(96)] 设工厂A和工厂B的产品次品率分别为$1\%$和$2\%$现产品中来自工厂A的概率是60\%,来自B的概率是40\%问,从这批
    商品中随机抽取一件,该件商品时次品且来自于A的概率是
    \item [(96)] 设$\xi,\eta$是两个相互独立且均服从正态分布$N(0,\frac{1}{2})$的随机变量,则随机变量$\left|\xi - \eta\right|$的期望是    
\end{enumerate}
\ifx\allfiles\undefined
\end{document}
\fi